\documentclass[11pt]{article}
\usepackage[margin=1in]{geometry}
\usepackage{graphicx}
\usepackage{amsmath,amssymb}
\usepackage{xcolor}
\usepackage{booktabs}
\usepackage[hidelinks]{hyperref}
\usepackage{parskip}

\title{\textbf{A true vertex model for epithelial mechanics in \texttt{plexus2}:\\
force balance, explicit T1, and what the Voronoi route could not do}}
\author{Notes from a \texttt{plexus2} prototype of \texttt{tyssue} (DamCB)}
\date{July 2026 --- \emph{living document, updated as the prototype progresses}}

\begin{document}
\maketitle

\begin{abstract}
\noindent
This is the companion to the Turing--vertex prototype. There we modelled epithelial mechanics
as a \emph{Self-Propelled Voronoi} (SPV) tissue --- each cell a point, the polygon its Voronoi
cell, the mesh re-tessellated every tick so neighbour exchanges (T1) were implicit. That report
closed on a clear diagnosis: clean \emph{tubulation} was gated not on chemistry but on two
mechanics ingredients the Voronoi route structurally lacks --- (i) \emph{force-balance iteration}
and (ii) an \emph{explicit reconnection} (T1) operator. Here we build the sibling implementation:
a \emph{true vertex model} (AVM) after \texttt{tyssue}, whose degrees of freedom are the
\emph{vertices} of a half-edge mesh, in which those two ingredients are native. The same
epithelial shape-energy contract, a genuinely different numerical realisation --- the ``one
contract, many implementations'' thesis of \texttt{plexus2}~\S5, made concrete. This document
records the model-deciphering as it happens; findings are added as each stage lands.
\textbf{[Results sections are filled in incrementally --- see the dated field notes.]}
\end{abstract}

\section{Two implementations of one contract}
Both prototypes realise the same biological operator --- an epithelial \emph{mechanical
interaction} whose energy is a per-cell area/perimeter shape energy
\[
E \;=\; \sum_f \Big[ K_A (A_f - A_f^0)^2 + K_P (P_f - P_f^0)^2 \Big]
        \;+\; \tfrac12\Gamma \sum_f P_f^2 \;+\; \Lambda \sum_e \ell_e/2 ,
\]
with a target shape index $p_0 = P^0/\sqrt{A^0}$ tuning a rigidity transition at
$p_0^\ast\!\approx\!3.81$. They differ only in \emph{how} the polygons and their motion are
represented:

\begin{center}
\renewcommand{\arraystretch}{1.3}
\begin{tabular}{@{}p{3.1cm}p{5.3cm}p{5.3cm}@{}}
\toprule
 & \textbf{SPV} (Turing\_vertex) & \textbf{AVM} (this prototype, \texttt{tyssue}) \\
\midrule
degrees of freedom & cell \emph{centres} (points) & mesh \emph{vertices} (half-edge mesh) \\
polygon & Voronoi cell of the centres & explicit face = ring of shared vertices \\
adjacency / T1 & re-tessellate each tick $\Rightarrow$ \emph{implicit} & \emph{explicit} local edge collapse\,+\,split \\
force balance & overdamped Euler steps & gradient descent to residual (quasistatic) \\
cost & Delaunay per tick (scipy) & one scatter-add + one autograd pass \\
\bottomrule
\end{tabular}
\end{center}

The AVM is the missing sibling: it supplies force balance and explicit T1, which is exactly
what the SPV report identified as the tubulation bottleneck. Concretely, both are registered
\texttt{plexus2} operators acting on a \texttt{set}; the SPV acts on a \texttt{cell} set, the
AVM on a \texttt{vertex} set with the faces carried as a half-edge table.

\paragraph{The original (Okuda) is a 3D vertex model --- this is the 2D rung of that ladder.}
Okuda \emph{et al.}\ is a \emph{true 3D} vertex model: cells are polyhedra sharing polygonal
faces, the DOF are the junction vertices $\mathbf r_i$ (overdamped
$\eta_i(\dot{\mathbf r}_i-\mathbf V_{fi})=-\nabla_i U$), the energy
$U=\sum_j \tfrac12 k_v(v_j-v_j^{\mathrm{eq}})^2+\kappa_s s_j$ is over each cell's polyhedral
\emph{volume} and \emph{surface}, the balance is iterated to residual $<10^{-5}$ every
$\Delta t_v$, neighbour changes are explicit \emph{reversible network reconnection} (RNR / T1)
every $\Delta t_r$, and division inserts a new face. The prototype here is the \emph{2D} rung of
that same ladder --- vertices, force balance, explicit T1 --- with an area/perimeter energy
standing in for volume/surface. That is deliberate: 2D is where T1 is a clean edge flip, so we
validate force balance and reconnection there first. tyssue itself carries the 3D rungs
(\texttt{Monolayer}/\texttt{Bulk}: \texttt{cell\_volume} elasticity, the IH/HI 3D reconnections,
\texttt{extrude} apical$\to$3D), all catalogued in the extraction (\S4) --- so the path to Okuda's
elongated \emph{tubes} (rather than the SPV prototype's differential-growth \emph{lobes}) is to
lift this 2D core to a monolayer once the 2D topology operators are validated.

\section{The model in one paragraph}
The mechanical set is the \texttt{vertex} (state \texttt{pos}); the tissue topology is a
\emph{half-edge table} $(\mathrm{srce},\mathrm{trgt},\mathrm{face})$ ordered counter-clockwise
around each face, stashed on the vertex \texttt{Level} exactly as the SPV prototype stashed its
Voronoi rings. Per-face targets $A^0,P^0$ and liveness live alongside it. \texttt{mesh\_seed}
bootstraps a clean honeycomb by taking the Voronoi tessellation of a triangular lattice
\emph{once} (outer rings dropped, their shared vertices pinned $\to$ a fixed-border patch), then
the AVM takes over. \texttt{shape\_energy} computes every face's area and perimeter with a
single \texttt{index\_add} over the half-edge table, forms the shape energy, and takes the force
$-\nabla E$ on the vertices by one autograd pass; an inner loop of \texttt{relax\_iters} gradient
steps per tick returns the net displacement as a velocity, so each rendered frame is
(near-)force-balanced --- the quasistatic force balance the original does with \texttt{scipy}
L-BFGS, here vectorised and differentiable. Boundary (pinned) vertices are held fixed.

\paragraph{A half-edge primer.}
A \emph{half-edge} is a \emph{directed} edge $A\!\to\!B$ that belongs to exactly one face; the
real edge shared by cells $L$ and $R$ is stored as two opposite half-edges. Each is three
integers $(\mathrm{srce},\mathrm{trgt},\mathrm{face})$:
\begin{center}\small
\begin{verbatim}
             cell L   (interior on the left)
     ... -> A --------------> B -> ...    half-edge  A->B ,  face = L
     ... <- A <-------------- B <- ...    half-edge  B->A ,  face = R
             cell R
   E_srce[k] = A     E_trgt[k] = B     E_face[k] = L   (one row per half-edge)
\end{verbatim}
\end{center}
A face is the ring of its half-edges in CCW order, so orientation --- and a positive shoelace
area --- is automatic; the two half-edges of a junction name its two neighbour cells, which is
exactly what a T1 needs (grab $A\!\to\!B$ in $L$ and $B\!\to\!A$ in $R$, collapse, re-split); and
per-face area/perimeter are a single \texttt{index\_add} keyed by \texttt{E\_face} (no per-cell
loop, so the mechanics vectorise). In \texttt{plexus2} terms \texttt{srce}/\texttt{trgt}/%
\texttt{face} are three index-buffer \emph{maps} on the set: \texttt{srce}/\texttt{trgt} are
edge$\to$vertex incidence, \texttt{face} is edge$\to$cell containment; gathering a cell value onto
its half-edges is a \emph{broadcast}, reducing half-edges to a cell is an \emph{aggregate}. Line
tension is $\Lambda\ell/2$ \emph{per half-edge}, so the two halves sum to $\Lambda\ell$ per real
edge --- tyssue's half-edge convention.

\section{Design decisions deciphered so far}
\emph{(These are the modelling choices that were not obvious from the paper and cost thought.)}

\paragraph{Vertices, not cells, are the set.}
The natural \texttt{plexus2} temptation is to keep the \texttt{cell} set of the SPV prototype and
bolt a mesh on. That is wrong: the AVM \emph{integrates the vertices}. Cells share vertices, so a
cell is not an independent particle --- it is a face, a derived readout over a ring of the vertex
set. The clean mapping is therefore a \texttt{vertex} set (the DOF) plus the half-edge table as
its topology map; faces never get their own integrated \texttt{Level}. Area/perimeter are
\emph{aggregate} readouts (vertices\,$\to$\,face), the force is their gradient back on the
vertices --- a genuine cross-scale (aggregate/broadcast) pair, richer than the single-set SPV.

\paragraph{Bootstrap the mesh from a Voronoi, then leave Voronoi behind.}
Hand-building a periodic honeycomb half-edge mesh with correct orientation is fiddly and
error-prone. Taking the Voronoi tessellation of a jittered triangular lattice \emph{once} gives a
valid honeycomb for free (scipy already de-duplicates the Voronoi vertices globally, so shared
vertices are shared indices). We keep only interior faces (finite Voronoi regions away from the
ragged border) and pin the boundary. After frame~0 the Voronoi is never used again --- the mesh
evolves by force balance and, later, explicit T1.

\paragraph{Force balance = inner relaxation loop, expressed as a velocity.}
A \texttt{plexus2} dynamics operator must \emph{return a delta}, not mutate positions. To emulate
tyssue's ``minimise to residual between topology events'' inside that contract, the operator runs
\texttt{relax\_iters} gradient-descent steps internally and returns
$v=(x_{\text{relaxed}}-x_0)/\Delta t$; the engine's $x \mathrel{+}= \Delta t\, v$ then lands on the
relaxed state. The $\Delta t$ cancels, so this is a pure per-frame relaxation and $\Delta t$ is a
free label --- the opposite of the SPV, where $\Delta t$ was the master knob that detuned the
mechanics. \emph{This alone answers the SPV report's ``force-balance iteration is missing''.}

\section{Stage 1: force-balanced relaxation and the rigidity transition}
A jittered honeycomb of $251$ interior cells ($566$ vertices, border pinned) is relaxed to
force balance by \texttt{shape\_energy} (\texttt{relax\_iters}$=20$ inner steps/frame,
$120$ frames), sweeping the target shape index $p_0\in[3.6,4.1]$. The residual (relaxed)
energy per cell and the relaxed mean shape index $\langle q\rangle = P/\sqrt{A}$ are the
order parameters of the rigidity transition (Bi \emph{et al.}, \emph{Nat.\ Phys.}\ 2015).

\begin{figure}[h]
\centering
\includegraphics[width=0.95\textwidth]{archive/_transition_tyssue2d.png}
\caption{\textbf{The AVM rigidity transition, from true force balance.} Left: residual energy
per cell after relaxation drops $\sim\!4$ orders of magnitude across $p_0^\ast\!\approx\!3.81$
(log scale) --- below it the tissue is \emph{rigid} (a finite energy barrier survives
relaxation), above it \emph{floppy} (zero-energy ground states exist). Right: the relaxed
$\langle q\rangle$ \emph{saturates} at the hexagonal floor $\approx\!3.72$ below the
transition (the cells cannot get rounder than a regular hexagon) and \emph{tracks} $p_0$ above
it (they hit their target shape at zero cost). No timestep tuning was required.}
\end{figure}

\begin{center}
\renewcommand{\arraystretch}{1.15}
\begin{tabular}{@{}rccl@{}}
\toprule
$p_0$ & residual $E$/cell & relaxed $\langle q\rangle$ & state \\
\midrule
3.60 & $1.50\times10^{-2}$ & 3.725 & rigid (barrier finite; $\langle q\rangle$ pinned at hex floor) \\
3.70 & $9.3\times10^{-4}$  & 3.725 & near-transition \\
3.75 & $8\times10^{-6}$    & 3.745 & \emph{knee} --- residual collapses here \\
3.78 & $2\times10^{-6}$    & 3.778 & floppy ($E\!\to\!0$; $\langle q\rangle$ tracks $p_0$) \\
3.81 & $1\times10^{-6}$    & 3.808 & floppy \\
3.84 & $1\times10^{-6}$    & 3.838 & fluid \\
3.90 & $1\times10^{-6}$    & 3.899 & fluid \\
4.00 & $2\times10^{-6}$    & 3.999 & fluid \\
4.10 & $2.1\times10^{-5}$  & 4.099 & fluid \\
\bottomrule
\end{tabular}
\end{center}

\begin{figure}[h]
\centering
\includegraphics[width=0.95\textwidth]{archive/solid_3p60/strip.png}\\[3pt]
\includegraphics[width=0.95\textwidth]{archive/fluid_4p10/strip.png}
\caption{\textbf{Rigid vs.\ floppy}, $0/15/50/100\%$ of relaxation. \emph{Top}
(\texttt{solid\_3p60}, $p_0\!=\!3.6$): the jittered honeycomb relaxes to a clean, jammed
hexagonal packing; every cell is blue (compact, shape index $<3.81$). \emph{Bottom}
(\texttt{fluid\_4p10}, $p_0\!=\!4.1$): the same tissue relaxes to elongated cells (red,
shape index tracking $p_0$) at essentially zero energy --- floppy. Border cells (pinned)
frame the bulk. (Without T1 the fluid tissue cannot yet \emph{flow} past its neighbours ---
that is Stage 2.)}
\end{figure}

\section{Operator extraction: the tyssue atlas}
Beyond completing the tubulation story, tyssue is a rich enough simulator to be worth
\emph{decomposing in full} into the Plexus operator algebra (plexus2 App.~B \& H). We read every
subsystem --- \texttt{dynamics}, \texttt{topology}, \texttt{geometry}/\texttt{core},
\texttt{behaviors}/\texttt{solvers}/\texttt{collisions} --- and normalised each mechanism to a
typed contract. The full record is \texttt{tyssue\_atlas.yaml}; the structure and the saturation
result are below.

\paragraph{One structural picture.} tyssue is a \emph{single half-edge \texttt{Level}}: only
\texttt{vert.pos} is an integrated mechanical DOF; the \texttt{edge} set carries every incidence
map ($\mathrm{srce},\mathrm{trgt},\mathrm{face},\mathrm{cell},\mathrm{opposite}$); \texttt{face}
and \texttt{cell} are pure aggregate readouts. Two engine primitives generate everything:
\texttt{upcast\_*} (gather a face/cell value onto its half-edges $=$ \textbf{broadcast}) and
\texttt{sum\_*}/\texttt{mean\_*} (scatter-reduce half-edges onto a vert/face/cell $=$
\textbf{aggregate}). Geometry (\texttt{update\_all}) is a feed-forward aggregate/broadcast DAG;
mechanics is an \emph{additive superposition} of effector forces onto one \texttt{vert.force}
accumulator; a \emph{solver} integrates it; an \emph{event manager} schedules structural changes.

\paragraph{On naming --- a deciphering note.} tyssue's central methods are
\texttt{compute\_energy}, \texttt{compute\_gradient}, \texttt{find\_energy\_min},
\texttt{update\_all}. \emph{None} is a Plexus operator. plexus2 names an operator by the biological
\emph{mechanism} it performs (the contract: \texttt{line\_tension}, \texttt{area\_elasticity},
\texttt{reconnect\_T1}, \texttt{divide\_cell}), and names the numerics separately as an
\texttt{implementation}. A generic numerical verb is therefore a \emph{signal} that one has found
not an operator but a \emph{composition} (\texttt{compute\_energy} $=$ the additive superposition
of the twelve mechanics contracts), a \emph{solver} (\texttt{find\_energy\_min} $=$ an
implementation of \texttt{force\_balance}), or a \emph{readout pipeline} (\texttt{update\_all} $=$
the geometry aggregate/broadcast DAG). The extraction's job is precisely to replace these verbs
with mechanism-named contracts.

\paragraph{Saturation.} The mechanics \emph{saturate} against the existing algebra --- the
shape-energy family is already in Plexus and tyssue merely adds (analytic-gradient)
implementations. The genuinely new contributions are in \emph{topology}, \emph{numerics} and
\emph{structure}. This is exactly the plexus2 prediction: a mature simulator contributes mostly
implementations of existing contracts, plus a few new operators where the language was thin.

\begin{center}
\footnotesize
\renewcommand{\arraystretch}{1.15}
\begin{tabular}{@{}p{2.6cm}p{7.6cm}p{2.4cm}@{}}
\toprule
\textbf{Subsystem} & \textbf{Normalised contracts} & \textbf{vs.\ Plexus} \\
\midrule
mechanics & line\_tension, perimeter\_elasticity, contractility, area\_elasticity,
volume\_elasticity, surface\_tension, lumen\_pressure & \textbf{existing} (saturated) \\
mechanics (new) & length\_elasticity, apico\_basal\_tension, boundary\_constraint,
edge\_friction & \textbf{new} \\
geometry & edge\_length, perimeter, face\_centroid, face\_area, cell\_volume, lumen\_volume
(aggregate/broadcast readouts) & existing \\
topology & divide\_cell / divide\_face, extrude\_face / remove\_cell & existing (divide / die) \\
topology (new) & \textbf{reconnect\_T1} (+3D IH/HI), mesh\_cleanup, repair\_pinched & \textbf{new} \\
numerics & overdamped\_euler & existing \\
numerics (new) & \textbf{force\_balance} (quasistatic minimise-to-residual), isotropic & \textbf{new} \\
structure (new) & \textbf{schedule} (EventManager), \textbf{collision\_constraint} & \textbf{new} \\
\bottomrule
\end{tabular}
\end{center}

The two boldfaced new contracts --- \texttt{reconnect\_T1} and \texttt{force\_balance} --- are
exactly the pair the Turing--vertex report named as the tubulation bottleneck. They are not
incidental: they are what a \emph{true} vertex model has and a Voronoi one cannot.

\section{Stage 2: explicit topology --- T1, division, extrusion}
The operators the Voronoi route could not have live in a forked module
(\texttt{tyssue\_topology\_ops.py}, so the mechanics and the topology can evolve separately):
\texttt{t1\_transition} (rewire; reversible network reconnection = collapse a sub-threshold edge
and split the merged vertex), \texttt{face\_divide} (structural; a septum splits a cell),
\texttt{face\_extrude} (structural; collapse a face to a point), and \texttt{face\_growth}
(structural; inflate the per-cell target area). They mutate the face-ring topology carried on the
vertex \texttt{Level} and rebuild the flat half-edge table that \texttt{shape\_energy} reads.

\paragraph{Robustness contract.} Every topology mutation is applied to a trial copy, the affected
faces are validated (three-plus unique vertices, positive signed area), and the change is
\emph{committed only if valid} --- otherwise it is a silent no-op. The mesh therefore never
tangles from a bad flip, which matters for unattended runs. A standalone self-test on the
honeycomb confirms it: T1 fired on $30/40$ shortest edges (the $10$ skips are boundary /
non-three-regular edges, correctly refused), four divisions and one extrusion all executed, and
\emph{zero} invalid faces resulted.

\paragraph{In-engine demonstration.} Running \texttt{mesh\_seed} $\to$ \texttt{face\_divide} $\to$
\texttt{shape\_energy} $\to$ \texttt{t1\_transition} through the engine, a one-shot clonal
division of $25\%$ of the cells followed by relaxation grows the tissue from $162$ to $198$ cells
and \emph{stays force-balanced with zero invalid faces} (Fig.~\ref{fig:clonal}). The division
operator is a genuine \texttt{plexus2} structural operator on the vertex set, and the sheet
re-relaxes around each new junction.

\begin{figure}[h]
\centering
\includegraphics[width=0.95\textwidth]{archive/clonal_divide/before_after.png}
\caption{\textbf{Clonal division on the force-balanced AVM} (\texttt{clonal\_divide}). Left: the
initial honeycomb ($F\!=\!162$). Right: after a one-shot division of $25\%$ of cells and
relaxation ($F\!=\!198$), coloured by shape index. The mesh stays valid (0 invalid faces) and
force-balanced --- the structural operator composes cleanly with the mechanics.}
\label{fig:clonal}
\end{figure}

\section{The two-level hierarchy: why \texttt{plexus2}, not just \texttt{tyssue}}
tyssue stores the whole tissue in one flat half-edge dataframe; a ``cell'' there is a \emph{row of
derived geometry}. \texttt{plexus2} instead makes the cell a first-class biological \emph{set}
carrying its \emph{own} state --- target area $a_0$, type/fate \texttt{ctype}, and (Goal~2)
morphogen \texttt{chem} --- none of which is derivable from vertex positions. The model then has
two Levels, \texttt{vertex} and \texttt{cell}, joined by the half-edge map, with a genuine
\texttt{cell\_geometry} \emph{Aggregate} (vertices\,$\to$\,cell area/perimeter) and a
\emph{cross-set} \texttt{shape\_energy} that reads $a_0$ from the cell set and broadcasts the force
back to the vertices (\texttt{INPUTS=[vertex, cell]}).

The benefit is concrete (Fig.~\ref{fig:clone}), and it must be done \emph{as biology, not as a
hack}. A first draft assigned fate by fiat --- a \texttt{cell\_paint} operator marked every cell
inside a hard-coded disc as a different type with a bigger target area. That proved the plumbing
but it is not a mechanism, so it was retired. In its place: a \texttt{cell\_morphogen} operator
imposes a smooth \emph{activator} field on the cell state \texttt{chem} (positional information ---
a Gaussian bump), and a \texttt{morphogen\_growth} operator makes each cell's target area a
\emph{continuous} response to it,
\[
a_0 \;=\; a_0^{\text{base}}\big(\rho + g\,\mathrm{Hill}(a)\big),\qquad
\mathrm{Hill}(a)=\frac{a^{n}}{a^{n}+a_{\text{sw}}^{n}} ,
\]
the Turing--vertex growth coupling, now acting on the cell set. The cross-set \texttt{shape\_energy}
reads $a_0$ and the tissue \emph{bulges where the activator is high}. The \emph{same} morphogen
field composes with mechanics in two ways, both principled: a \emph{continuous} response
(\texttt{morphogen\_growth}, above) gives a graded bulge whose cell area tracks the activator with
correlation $0.97$; and a \emph{discrete} response (\texttt{cell\_differentiate}) thresholds the
activator into a genuine cell \emph{type} --- a French-flag fate --- so the cells above threshold
become a coherent \emph{clone} (Fig.~\ref{fig:clone}) with a larger target area that bulges to
$2.6\times$ the wild-type, all force-balanced. In tyssue either coupling would be bespoke glue; here
each is \emph{state on a set} plus a typed operator (\texttt{morphogen $\to$ growth}, or
\texttt{morphogen $\to$ threshold $\to$ typed mechanics}), and it is exactly the substrate Goal~2
needs --- swap the imposed bump for a live Turing reaction--diffusion on the cell--cell adjacency and
the clone becomes a \emph{self-organised, activator-driven bud}. That composability (mechanism $+$
mechanism through a shared set/operator language) is the payoff of decomposing tyssue into
\texttt{plexus2} rather than wrapping it.

\begin{figure}[h]
\centering
\includegraphics[width=0.98\textwidth]{archive/clone_bulge/strip.png}
\caption{\textbf{A genuine cell set: a morphogen differentiates a typed clone that drives mechanics}
(\texttt{clone\_bulge}), $0/33/66/100\%$. \texttt{cell\_morphogen} imposes a smooth activator field;
\texttt{cell\_differentiate} thresholds it into a fate (a French-flag rule), so the cells above
threshold become a coherent \emph{clone} (red) of a genuine cell type with a larger target area; the
cross-set \texttt{shape\_energy} makes that clone bulge to $2.6\times$ the wild-type (white) cell
area. The clone is \emph{derived from the morphogen}, not a hard-coded disc --- the retired
\texttt{cell\_paint} assigned fate by fiat; here fate is state on a set, set by an operator that
reads a chemical field. This is the bridge to Goal~2 (replace the imposed bump with a live
reaction--diffusion and the clone self-organises).}
\label{fig:clone}
\end{figure}

\section{Recapitulating the \texttt{tyssue-demo} notebooks}
The canonical \texttt{tyssue} tutorials (\texttt{DamCB/tyssue-demo}) are a natural checklist: one
demo per mechanism. Rebuilding each as a \texttt{plexus2} operator composition is precisely
Goal~1. Current coverage ($\checkmark$ sound, $\sim$ partial, $\rightarrow$ next):

\begin{center}
\small
\renewcommand{\arraystretch}{1.15}
\begin{tabular}{@{}p{3.9cm}p{4.7cm}p{3.5cm}c@{}}
\toprule
\textbf{notebook} & \textbf{shows} & \textbf{our demo} & \\
\midrule
00-Basic\_Usage       & build / inspect a sheet          & \texttt{mesh\_seed} + cell set        & $\checkmark$ \\
01-Geometry           & area / perimeter / length        & \texttt{cell\_geometry} (aggregate)   & $\checkmark$ \\
02-Visualization      & draw the polygons                & strip + movie renderers               & $\checkmark$ \\
03-Dynamics           & effectors / energy               & \texttt{shape\_energy}                & $\checkmark$ \\
04-Solvers            & QS / Euler solvers               & bounded differentiable Euler          & $\checkmark$ \\
D-FarhadifarModel     & the AVM energy \& rigidity        & \texttt{run\_tyssue2d} transition     & $\checkmark$ \\
05-Rearangements      & T1 neighbour swaps               & \texttt{t1\_transition} + flow        & $\checkmark$ \\
C-2DMigration         & active migration                 & \texttt{run\_tyssue\_flow}            & $\checkmark$ \\
07-EventManager       & scheduling behaviours            & the \texttt{plexus2} schedule         & $\checkmark$ \\
06-Cell\_Division     & straight-line division by a plane & \texttt{face\_divide\_line}          & $\checkmark$ \\
B-Apoptosis           & cell death / extrusion           & \texttt{apoptosis}                    & $\checkmark$ \\
A-RNR (2D)            & multistep reconnection           & \texttt{apoptosis} (shed\,$\to$\,collapse) & $\checkmark$ \\
\bottomrule
\end{tabular}
\end{center}

All eleven 2D-relevant notebooks now have a sound \texttt{plexus2} recapitulation. What remains is
not a 2D mechanism but the \emph{lift to 3D}: the monolayer / bulk geometry and the 3D IH/HI
reconnections (the 3D half of A-RNR), which are the 3D-vesicle stage (Goal 1bis). The 2D operator
set --- mechanics, force balance, T1, division, extrusion, RNR, and the two-level morphogen
coupling --- is complete and cross-validated against the tyssue tutorials.

\paragraph{Straight-line (oriented) division (06).}
The placeholder through-vertex split (which degraded after $1$--$2$ rounds) is replaced by the
proper Brodland--Veldhuis division at a \emph{decidable angle}: the plane through a cell's centroid
cuts its two crossed edges, a new vertex is inserted on each (shared with the neighbour across that
edge), and a septum joins them --- well-shaped daughters, each targeting \emph{half} the mother's
area (area conserved, as tyssue halves \texttt{prefered\_area}). Fig.~\ref{fig:divide} divides
$35\%$ of a relaxed sheet at a fixed angle $\pi/2$ and re-relaxes: $118\!\to\!159$ cells, $0$ invalid
faces, the oriented (vertical) septa clearly visible --- the decidable division plane the demo is
about.

\begin{figure}[h]
\centering
\includegraphics[width=0.98\textwidth]{archive/divide_pi2/strip.png}
\caption{\textbf{Straight-line division at a decidable angle} (\texttt{divide\_pi2}), coloured by
shape index (blue solid / red fluid). A relaxed honeycomb (left) divides $35\%$ of its cells by a
plane at $\pi/2$ (vertical septa), adding a vertex on each crossed edge, then re-relaxes to
$159$ cells with $0$ invalid faces. Recapitulates \texttt{tyssue-demo} 06-Cell\_Division.}
\label{fig:divide}
\end{figure}

\paragraph{Apoptosis / cell elimination (B) + reversible network reconnection (A-RNR).}
The \texttt{apoptosis} operator shrinks the target area of a chosen cell each tick (so
\texttt{shape\_energy} contracts it) and, once it is small, eliminates it. A naive one-step
extrusion --- collapse the whole cell to a point --- leaves its six neighbours on a single vertex,
a degree-$6$ \emph{rosette} the three-regular T1 cannot relax (and collapsing a near-zero sliver is
numerically unstable). The fix is \emph{multistep} reconnection (Okuda's RNR; the A-RNR notebook, in
2D): while the cell is small it \emph{sheds a side} by a T1 on its shortest edge
($6\!\to\!5\!\to\!4\!\to\!3$), and only the resulting \emph{triangle} is collapsed --- to a clean
degree-$3$ junction. Cell (cortical) \emph{contractility} $\Gamma$ and line tension $\Lambda$ keep
the shrinking cell and its neighbours round rather than spiky (the neighbours \emph{constrict in},
apically closing the gap). With that, the cell is eliminated ($56\!\to\!55$, $0$ invalid faces) and
the tissue \emph{heals}: the neighbours reconnect and relax back to epithelium, no rosette
(Fig.~\ref{fig:apop}, the dying cell drawn dark grey); the only residue is the local topological defect
inherent to removing a cell from a hexagonal packing. This closes B-Apoptosis \emph{and} the 2D case
of A-RNR; the 3D IH/HI reconnections belong to the 3D-vesicle stage.

\begin{figure}[h]
\centering
\includegraphics[width=0.98\textwidth]{archive/apoptose/strip.png}
\caption{\textbf{Single-cell apoptosis with reversible network reconnection} (\texttt{apoptose}),
coloured by shape index. A central cell shrinks and is eliminated by \emph{multistep} reconnection:
it sheds sides by T1 until it is a triangle, which collapses to a degree-$3$ junction. The cell is
removed ($118\!\to\!117$, $0$ invalid), the neighbours reconnect, and the site heals back to
force-balanced epithelium --- no rosette, no residual distortion. Recapitulates \texttt{tyssue-demo}
B-Apoptosis and the 2D case of A-RNR.}
\label{fig:apop}
\end{figure}

\paragraph{Integration test: a growing, dividing tissue (and why surface tension matters).}
The operators compose into a developing tissue: \texttt{face\_growth} inflates each cell's target
area, \texttt{face\_divide\_line} in cell-cycle mode divides any cell past a threshold at a random
angle (edge-midpoint, well-shaped daughters each targeting half), \texttt{shape\_energy} holds force
balance, and \texttt{t1\_transition} resolves short junctions. On a \emph{free} (unpinned) boundary
the sheet expands as it grows. Three settings proved essential. (i)~A free boundary, so growth is not
frustrated into over-compression. (ii)~\emph{Line tension} $\Lambda\sum_e\ell_e$ --- without which the
dividing cells relax into sharp triangular slivers and the free rim spikes into points; it smooths the
boundary and straightens junctions. (iii)~\emph{Cortical contractility} $\tfrac12\Gamma\sum_f P_f^2$
--- a per-cell perimeter penalty that rounds each daughter out between divisions instead of letting it
elongate; without it the tissue drifts uniformly fluid (shape index far above $p_0$) and cells
sharpen. With line tension \emph{and} contractility the daughters stay round, well-shaped polygons at
a moderate shape index (this pair also smooths the apoptosis reconnection). The tissue proliferates
$88\!\to\!176$ cells with $0$ invalid faces --- \emph{sustained} edge-midpoint division
(Fig.~\ref{fig:grow}), in contrast to the Stage-2 through-vertex placeholder that tangled after
$1$--$2$ rounds.

\begin{figure}[h]
\centering
\includegraphics[width=0.98\textwidth]{archive/grow_divide/strip.png}
\caption{\textbf{A growing, dividing tissue} (\texttt{grow\_divide}), coloured by shape index. A free
honeycomb grows (target areas inflate) and cells divide by straight lines at random angles when they
pass a size threshold; the sheet expands and re-tessellates, staying valid ($88\!\to\!176$ cells,
$0$ invalid). Line tension ($\Lambda$) plus cortical contractility ($\Gamma$) keep the daughters
round and the free rim smooth rather than sharp/triangular. This is the sustained, well-shaped
proliferation the Stage-2 placeholder could not do.}
\label{fig:grow}
\end{figure}

\section{Goal 1bis: the 3D vesicle (surface vertex model)}
The same operators lift to 3D. A closed spherical half-edge mesh --- the spherical Voronoi of a
Fibonacci sphere, a Goldberg polyhedron of hexagons and exactly twelve pentagons, with Euler
characteristic $V\!-\!E\!+\!F=2$ and every edge shared by two faces --- carries the AVM in three
dimensions. \texttt{shape\_energy\_3d} uses the 3D geometry (a Newell area vector per face, 3D edge
lengths) plus a \emph{lumen-volume} term $K_V(V-V_0)^2$ that keeps the shell inflated against surface
tension --- the closure a 2D patch got from pinning. The force is one 3D autograd pass; the bounded
Euler step carries over unchanged.

One 3D-specific lesson earned its keep. A free surface with area $+$ tension energy and \emph{no
bending rigidity} \emph{buckles} out of plane (roughness $0.004\!\to\!0.155$) --- exactly as a real
epithelium does; the lumen term fixes the total volume but not local smoothness. Adding a soft
radial term $K_R\sum_v(|\mathbf r_v|-R_0)^2$ (a bending-like penalty toward the shell) recovers a
smooth relaxed vesicle: the roughness stays $\sim\!0.02$ and the cross-section stays circular
(Fig.~\ref{fig:ves}). This closed, force-balanced shell is the substrate for 3D morphogen-driven
budding --- the 3D IH/HI reconnections (the 3D half of A-RNR) and a mechanical monolayer (apical/basal,
with \texttt{cell\_volume} elasticity) are the remaining 3D increments.

The vesicle is rendered as a \emph{monolayer}: each cell is drawn as a prism --- an apical face
(outer), a basal face (inner, the surface scaled radially inward), and the lateral walls between ---
so a cell carries many faces, not one flat polygon. The apical faces share the mesh vertices and tile
edge-to-edge into a closed opaque shell; cells take the Turing white$\to$red activation LUT (here
activation $=0$, all white, awaiting the Goal 2 reaction--diffusion). The cross-section slices the
monolayer: each junction edge crossing the plane gives an apical point and its inward-scaled basal
counterpart, and the connecting segment is that junction's lateral wall, partitioning the shell into
the individual cells around the lumen (Fig.~\ref{fig:ves}, bottom).

\begin{figure}[h]
\centering
\includegraphics[width=0.98\textwidth]{archive/vesicle/strip.png}
\caption{\textbf{A 3D epithelial vesicle} (\texttt{vesicle\_3d}), $0/33/66/100\%$. \emph{Top}: the 3D
monolayer shell --- cells as opaque, edge-to-edge prisms (apical/basal/lateral faces), white
(activation $=0$) with black junctions. \emph{Bottom}: the monolayer cross-section, rendered to the
same scale --- a ring of individual cells (apical outer, basal inner, lateral walls) around the hollow
lumen. A jittered sphere relaxes to a uniform, smooth epithelial shell, held inflated by the volume
term and kept from buckling by the soft radial term. The true-vertex-model sibling of the
Turing\_vertex spherical-Voronoi vesicle.}
\label{fig:ves}
\end{figure}

\subsection*{Growth by per-cell volume elasticity (the emergent inflation)}
Growing the vesicle turned out to be the sharpest model-deciphering lesson of the 3D stage. The naive
recipe --- ramp each cell's target area $A_0$ and a single \emph{global} lumen volume $V_0$, then let
the shape energy inflate the shell --- \emph{fails}: the sphere lags its targets and the surplus cell
area buckles out of plane into deep crumples (roughness $0.01\!\to\!0.2$, shape index climbing from
$3.8$ to $4.4$ --- ``too much fabric on a ball''). Neither slowing the growth, adding relaxation, nor
strengthening a global lumen pump fixed it, because a \emph{single scalar} $K_V(V-V_0)^2$ constrains
only the \emph{sum} of the cell volumes: any set of local dips that cancel in the sum is energetically
free, so buckling is unopposed.

The fix is the mechanism both references already use. tyssue's \texttt{ClosedMonolayer} gives every
\emph{cell} a \texttt{prefered\_vol} and \texttt{vol\_elasticity} (not just a lumen), and
Turing\_vertex's Eq.~3 writes the cell energy as $U=\sum_j\big[\tfrac12 k_v(v_j-v_{\mathrm{eq},j})^2
+\kappa_s s_j\big]$ --- a \emph{per-cell} volume elastic energy. So we split the one global lumen term
into a sum over cells: the lumen volume $V=\sum_f v_f$ with per-cell wedge volumes
$v_f=\tfrac13(\mathbf{c}_f\!\cdot\!\mathbf N_f)$ (the pyramid from the sphere centre to face $f$), and
the energy carries $\sum_f K_V\,(v_f-v_{\mathrm{eq},f})^2$. Now each cell holds \emph{its own} volume:
a local dip lowers that cell's $v_f$ and is penalised, so the term \emph{actively resists buckling},
and growth is just ramping each $v_{\mathrm{eq},f}$ (with $A_0\propto s^2$, $v_{\mathrm{eq}}\propto
s^3$). The whole shell then inflates purely by force balance --- no vertex is moved by hand --- from
radius $\sim\!5$ to $\sim\!7.1$ while staying smooth (roughness $\le 0.01$) and keeping the cells at
their target shape (index $3.8\!\to\!3.87$), Fig.~\ref{fig:vesgrow}. The lesson: \emph{distributed}
elastic targets, not a global constraint, are what make morphogenetic expansion emerge and stay
stable --- the same reason real epithelia use cell-autonomous volume control.

\begin{figure}[h]
\centering
\includegraphics[width=0.98\textwidth]{archive/vesicle_grow/strip.png}
\caption{\textbf{A growing vesicle} (\texttt{vesicle\_grow}), $0/33/66/100\%$, framed to the final
extent. Ramping each cell's target volume $v_{\mathrm{eq}}$ inflates the shell through the per-cell
volume elasticity alone: the sphere (top) and its cross-section ring (bottom) grow self-similarly from
$R\!\approx\!5$ to $\approx\!7.1$, the cells stay uniform and well-tiled, and the shell stays smooth
(roughness $\le0.01$). Replacing this per-cell term with a single global lumen volume instead crumples
the shell --- the emergence of stable expansion requires distributed, not global, volume targets.}
\label{fig:vesgrow}
\end{figure}

\subsection*{Division in 3D (a proliferating vesicle)}
The last 3D increment is cell \emph{division} on the closed shell --- the sphere analog of tyssue's
\texttt{sheet\_topology.cell\_division} (a face split by a line; the bulk \texttt{monolayer\_topology}
version with a \texttt{vertical}/\texttt{horizontal} orientation waits for the mechanical monolayer).
The cell cycle is \emph{volume doubling}: each cell records its birth wedge volume and divides once its
current volume reaches twice that (\texttt{divide\_3d}). A division splits the cell by an
\emph{edge-midpoint septum} --- two of its edges are cut at their midpoints ($m_1,m_2$) and a new edge
$m_1$--$m_2$ separates the daughters. The subtlety unique to 3D: every edge of a closed surface is
shared by two faces, so cutting an edge \emph{also} inserts the midpoint into the neighbouring cell's
ring. Getting that bookkeeping right keeps the mesh a valid closed surface: each division adds exactly
$\Delta V\!=\!2,\ \Delta E\!=\!3,\ \Delta F\!=\!1$, so $V\!-\!E\!+\!F$ is invariant. A standalone check
(reorienting the SphericalVoronoi outward first, so every directed edge is unique) confirms
Euler $=2$ across dozens of divisions; \emph{in the engine} the growing, dividing vesicle proliferates
$150\!\to\!521$ cells over $371$ divisions in $220$ frames and still verifies \texttt{closed = true,
Euler = 2} (Fig.~\ref{fig:vesdiv}). Pacing follows the Turing\_vertex \texttt{fig4\_vesicle} reference
(volume-doubling, a few divisions per frame) so the tissue proliferates gradually. Two subtleties were
worth fixing: divisions must be chosen \emph{unbiased} --- selecting the first ready cells in face-index
order sweeps pole-to-pole because the Fibonacci seeding numbers faces that way, so we shuffle the
candidates and randomise the septum direction --- and daughters must be kept \emph{compact}. On the
last point we followed a literature sweep (tyssue's \texttt{monolayer\_topology}, Turing\_vertex, cellGPU):
the emergent fix is not a smoothing pass but two mechanisms already in the physics. (i)~\emph{Long-axis}
(Hertwig) division --- the septum cuts perpendicular to the cell's longest axis (an SVD of the cell's
vertices), which yields compact daughters instead of elongated slivers. (ii)~\emph{Cortical
contractility} $\tfrac12\Gamma\sum_f P_f^2$, the perimeter term that rounds each cell. Together they
drop the shape index from $\sim\!4.2$ to $\sim\!3.9$ \emph{emergently} (a geometric Lloyd smoothing
would also round the cells, but it is a fairing hack, not physics, so we discarded it). The last step
toward the regular-hexagon floor is \emph{T1 neighbour exchange}, the annealing mechanism, now added as
\texttt{reconnect\_t1\_3d} (a \textbf{Rewire} op, the 3D sibling of the 2D \texttt{t1\_transition}):
every interior edge shorter than a threshold is flipped by a local four-cell neighbour exchange
($A$--$B$ lose the edge, the far cells $C$--$D$ gain it), collapsing the edge to its midpoint and
reopening it perpendicular in the tangent plane. It keeps $V,E,F$ (and Euler) fixed and only reconnects;
a standalone check confirms the shell stays closed with Euler$=2$ through $200+$ flips. Adding it to the
dividing vesicle anneals the division defects and drops the shape index a further $3.96\!\to\!3.88$.
Note the sphere mandates $\ge\!12$ pentagonal defects (Euler), so perfect hexagons are not the target. \emph{At scale} (the fig4 test: $600\!\to\!2367$ cells over $1000$ frames) the shell stays
\texttt{closed, Euler}$=2$, roughness $0.008$, shape index $3.84$ --- rounder at scale, as more cells
anneal better.

\begin{figure}[h]
\centering
\includegraphics[width=0.98\textwidth]{archive/vesicle_grow_divide/strip.png}
\caption{\textbf{A dividing (proliferating) vesicle} (\texttt{vesicle\_divide}), $0/33/66/100\%$,
framed to the final extent. The vesicle grows by per-cell volume elasticity \emph{and} its cells divide
by a volume-doubling cell cycle (edge-midpoint septum): the shell inflates from $R\!\approx\!5$ to
$\approx\!9.5$ while the cell count grows $150\!\to\!521$, so cells stay near constant size
(top: 3D shell; bottom: cross-section, a ring of ever more, smaller cells around the lumen). Division
is chosen unbiased (shuffled candidates, random septum direction) so it fills the sphere uniformly. The
closed surface stays valid throughout ($\texttt{Euler}=2$ after all $371$ divisions), line tension
keeping the fresh daughters well-shaped.}
\label{fig:vesdiv}
\end{figure}

\section{Goal 2: Turing reaction--diffusion on the cell set}
The two-level hierarchy pays off again for morphogenesis. The morphogen lives on the \emph{cell} set as
state \texttt{chem}$=[a,h]$ (a first-class biological DOF, not derivable from vertex geometry), while
the mechanical \emph{vertex} set holds the shell. Three operators run the reaction--diffusion each
frame: \texttt{cell\_geometry\_3d} (an \textbf{Aggregate}: vertices $\to$ per-cell centroid),
\texttt{cell\_adjacency} (a \textbf{Rewire}: the cell--cell neighbour graph) and
\texttt{cell\_diffuse}\,$+$\,\texttt{cell\_react} (\textbf{Lateral} field updates that integrate
\texttt{chem}).

Two points are worth stating. First, \emph{how we get cell neighbours without a Voronoi}: in the AVM
the cells \emph{are} the faces of the half-edge mesh, so two cells are neighbours iff they share a mesh
\emph{edge} --- \texttt{cell\_adjacency} groups the half-edges by their vertex pair and links the two
faces that share each edge. The mesh topology already \emph{is} the adjacency graph; no Delaunay/Voronoi
of cell centres is needed (that is only how the Self-Propelled-Voronoi sibling builds \emph{its} graph).
Second, \texttt{react} is a plexus2 \emph{contract} with interchangeable implementations: on the same
cell set, \texttt{gray\_scott} gives a coral/labyrinth and \texttt{brusselator} (its parameters
transposed verbatim from Turing\_vertex \texttt{fig4\_coral}: $\gamma=2,A=1,B=3$; $d_a=0.05, d_h=0.7$)
gives round Turing spots. The Brusselator's fast-inhibitor diffusion is only stable at the fig4
timestep $\mathrm{d}t=0.02$, so the whole sim runs at that step and \texttt{shape\_energy\_3d} is given
the same $\mathrm{d}t$ so the mechanics still relaxes fully each frame. A self-organised activator
pattern emerges across the closed shell (Fig.~\ref{fig:coral}), coloured white$\to$red.

Crucially the pattern couples to \emph{topological change}. By construction \texttt{rd\_coral\_grow}
is \emph{exactly the dividing vesicle of Fig.~\ref{fig:vesdiv} plus the coral colouring}: its mechanical
schedule (growth $\to$ shape energy $\to$ T1 $\to$ division) is that of \texttt{vesicle\_divide}
verbatim, and the reaction--diffusion ops only \emph{read} the mesh (aggregate vertices to cell
centroids, integrate \texttt{chem}) --- they never move a vertex. We verified that the hard way: a
controlled triple at identical mechanics --- \texttt{grow\_only}, \texttt{grow\_divide}, and
\texttt{grow\_divide}$+$RD --- gave \emph{byte-identical vertex trajectories} for the last two, so the
morphogen is provably pure colouring. The one genuine coupling runs the other way:
\texttt{divide\_3d} propagates the mother's \texttt{chem} to its daughters (the kept face retains the
mother's state, the new face inherits a copy), so the Turing field rides the proliferating tissue
coherently instead of being scrambled by division. At the vesicle\_divide scale ($150\!\to\!518$ cells)
the labyrinth develops cleanly on the growing, dividing shell.

\paragraph{The synchronised division wave, and the stochastic cell cycle that breaks it.} Scaling the
same recipe up --- to $500$ cells over $1000$ frames, for a finer pattern --- exposed a real
instability, and it took the controlled experiment above to show it was \emph{not} the
reaction--diffusion. Because every cell is born at the same volume and grows uniformly, they all cross
the volume-doubling threshold \emph{in lockstep}: the cell count climbs in flat $+100$-per-frame steps
(a wave capped only by \texttt{max\_div}), so the shell is driven under sustained \emph{max-rate}
division for the entire growth ramp. Fresh, unrelaxed daughters accumulate faster than the bounded
Euler relaxation can smooth them and the (face-referenced) shell roughness at mid-ramp blows up to
$0.38$ (worst frame $1.7$) --- the transient ``grey'' of exposed cell walls --- annealing back only
once growth plateaus. The fix is the one real epithelia use: \emph{cell-cycle-time variability}. Two
knobs make the division \emph{stochastic}: a \emph{stochastic volume seed} (\texttt{vseed\_cv} on
\texttt{seed\_mesh\_3d}) gives each cell a random initial cell-cycle phase so the \emph{first} wave is
desynchronised, and a \emph{stochastic cell cycle} (\texttt{cycle\_cv} on \texttt{divide\_3d}) draws
each daughter's next division threshold from a Gaussian so later generations stay spread out. With
CV$=0.4$ the lockstep wave becomes a smooth ramp (cells climb gradually from frame $0$), the mid-ramp
roughness drops two orders of magnitude ($0.38\!\to\!0.004$), and \texttt{rd\_coral\_grow\_big}
proliferates $500\!\to\!1994$ cells while staying smooth throughout --- a finer coral labyrinth on the
larger shell (Fig.~\ref{fig:coralgrow}). The lesson repeats the document's refrain: the artefact
\emph{looked} like a shading or reaction--diffusion bug, but it was a \emph{synchronisation} property
of the mechanics, cured by restoring the biological noise a deterministic cell cycle had omitted.

\begin{figure}[h]
\centering
\IfFileExists{archive/rd_coral_grow_big/strip.png}%
  {\includegraphics[width=0.98\textwidth]{archive/rd_coral_grow_big/strip.png}}%
  {\includegraphics[width=0.98\textwidth]{archive/rd_coral_grow/strip.png}}
\caption{\textbf{Coral reaction--diffusion riding a growing, dividing vesicle}
(\texttt{rd\_coral\_grow\_big}), $0/33/66/100\%$. The Gray-Scott labyrinth (cell state \texttt{chem},
coloured white$\to$red by activator) develops on a shell that grows by per-cell volume elasticity and
proliferates $500\!\to\!1994$ cells; \texttt{divide\_3d} hands each daughter a copy of the mother's
morphogen so the pattern stays coherent across division. A \emph{stochastic cell cycle} (per-cell
random cycle length, CV$=0.4$) desynchronises division: without it the cells double in lockstep and the
shell roughens $\sim\!100\times$ at mid-growth (a grey band of exposed walls); with it the shell stays
smooth throughout (top: 3D; bottom: cross-section).}
\label{fig:coralgrow}
\end{figure}

\paragraph{The ``hollow cells'' were a recording artefact, not a physical defect (2026-07-21).} A long
\texttt{rd\_coral\_grow} (and any dividing vesicle past $\sim\!500$ cells) rendered a shell the
\texttt{hollow\_flags} diagnostic scored as $\sim\!97\%$ ``hollow'' (folded/inverted caps) --- a striking
grey/red crumpling that annealed away by the final frame. A long chain of controlled experiments
established, in order, that it was division-driven, that a \emph{single} division is geometrically clean
(the Hertwig septum midpoints land \emph{on} the surface and both daughters' normals stay parallel to the
mother, $\mathbf{n}\!\cdot\!\mathbf{n}_{\text{mother}}=0.995,\,0.998$), and that \emph{no} physical remedy
touched it --- an anti-inversion filtered step, local daughter relaxation, dihedral bending to
$K_{\text{bend}}=100$, removing the radial term, tripling the relaxation, even slowing growth (which was
\emph{worse}). That every plausible mechanical fix failed was the tell: there was no mechanical defect.

The cause was a \textbf{recording misalignment}. The engine records vertex \emph{positions}
(\texttt{posf}) at one stride while \texttt{topo\_snapshot\_3d} records the changing \emph{mesh topology}
(\texttt{hist}) at another; on runs long enough to trip a recording stride $>1$ the two arrays end up
\emph{different lengths} (at $300$ frames \texttt{posf} has $151$ entries, \texttt{hist} $76$). The
diagnostic and renderer then paired positions from one frame with the half-edge connectivity of
\emph{another} --- ring vertex indices no longer matched the coordinates, every Newell normal was computed
from a scrambled polygon, and the metric reported phantom inversions. The proof is one line: the identical
$300$-frame run with \texttt{topo\_snapshot every}$=2$ (\texttt{posf}$=151$, \texttt{hist}$=76$,
\emph{misaligned}) scores hollow $0.98$; with \texttt{every}$=1$ (both $301$, \emph{aligned}) it scores
$\mathbf{0.000}$ throughout. The $220$-frame vesicle looked clean only because at that length the stride was
$1$ and the recordings happened to stay aligned. Recording topology every frame the positions are recorded
makes \texttt{vesicle\_grow\_divide} a smooth sphere all the way to $949$ cells (Fig.~\ref{fig:hollow}).
The lesson is the document's oldest refrain, sharpened: \emph{measure the geometric thing you mean} --- a
diagnostic that silently pairs two arrays of different provenance will fabricate a defect that no amount of
physics can fix. (One genuine, mild residue survives the fix: with a deterministic cell cycle the cells
divide in backlogged waves, so their \emph{areas} spread; the stochastic cell cycle CV$=0.4$ plus a
live-scaled division cap even them out, as \texttt{rd\_coral\_grow\_big} already showed. The earlier
\emph{roughness}-metric ``synchronised wave'' result is unaffected --- that metric reads positions only.)

\begin{figure}[h]
\centering
\includegraphics[width=0.98\textwidth]{archive/vesicle_grow_divide_400f/strip.png}
\caption{\textbf{With aligned recording the dividing vesicle is a smooth sphere --- the ``hollow
buckling'' was a diagnostic artefact.} The same physics that scored $\sim\!97\%$ ``hollow'' under
mismatched position/topology recording strides (\texttt{posf}$=151$, \texttt{hist}$=76$ at $300$ frames)
scores hollow $\le\!0.003$ once the two are recorded on the same stride
(\texttt{vesicle\_grow\_divide\_400f}, $150\!\to\!949$ cells; top row white render, bottom row coloured by
the hollow flag). The ``defect'' lived entirely in the diagnostic's frame pairing; the one real residue is
a mild cell-\emph{area} spread from the deterministic cell cycle.}
\label{fig:hollow}
\end{figure}

\paragraph{Cell-size homogenisation --- a hypothesis-driven fix (2026-07-22).} On the growing, dividing
vesicle the cells were too non-uniform (per-cell area CV $\approx\!0.6$; the tube-tip cell volume was
$1.5\times$ the base with a huge spread, $1.56\pm2.45$ vs $1.07\pm0.35$). Rather than a grid search, three
falsifiable hypotheses were tested one factor at a time (short $150$-cell/$180$-frame A/Bs, metric $=$
per-cell area CV). \emph{H1 --- grow by proliferation, not inflation:} \texttt{vesicle\_growth} enlarges
the shell by ramping every cell's target volume (cells bulge), whereas Okuda grows by \emph{proliferation}
--- capping $v_{eq}$ at $\tfrac43 v_{ref}$ so each cell cycles in $[\tfrac23,\tfrac43]v_{ref}$ and the
sphere grows by cell \emph{count}. Switching to \texttt{morphogen\_growth\_3d} with baseline $\rho=1$ and
the cap cut area CV $0.63\!\to\!0.40$ (validated). \emph{H2 --- tighten division timing:} the dominant
residue was the loose division-threshold jitter (CV$_{\text{cycle}}=0.4$); tightening to $0.15$ and
\emph{bounding the cell-cycle duration} (\texttt{min\_cycle}$=4$, \texttt{max\_cycle}$=12$) dropped area
CV $0.40\!\to\!0.18$ --- the biggest single step. \emph{H3 --- stiffen volume elasticity:} raising
$K_V\,1\!\to\!4$ so cells track $v_{eq}$ took area CV to $\mathbf{0.17}$ (Okuda-uniform range) and the
hollow fraction to $\mathbf{0}$: a spotless, uniform shell. The division model is now Okuda's ---
\emph{volume-primary} ($2\times$ birth volume) \emph{with a bounded cell-cycle duration} --- the two
competing timing models parameterised together. This uniform substrate is the prerequisite for the
multi-tube morphogenesis, and it is reached by short, hypothesis-driven steps rather than a blind sweep.

This is the substrate for the morphogenesis proper --- morphogen-driven growth (bulge/divide where $a$
is high) on a \emph{fluid} ($p_0\!>\!3.81$) shell where T1 lets the tissue flow, the bridge to tubes.

\begin{figure}[h]
\centering
\includegraphics[width=0.98\textwidth]{archive/coral/strip.png}
\caption{\textbf{Live Turing reaction--diffusion on the vesicle cell set} (\texttt{coral}, Gray-Scott),
$0/33/66/100\%$. The morphogen is CELL state; diffusion runs on the cell--cell adjacency read straight
from the half-edge table (no Voronoi). From faint noise a self-organised coral/labyrinth activator
pattern grows across the closed shell (top: 3D, coloured white$\to$red by activator; bottom:
cross-section). The \texttt{spots} preset swaps in the Brusselator (fig4 params) on the same set for
round Turing spots --- \texttt{react} is an interchangeable contract.}
\label{fig:coral}
\end{figure}

\section*{Field notes}
\emph{Dated log of what actually worked or failed --- the hard-won part.}

\paragraph{3D division at scale (2026-07-21) --- a synchronised division wave, not a rendering bug.}
\begin{itemize}\itemsep2pt
\item \textbf{The controlled experiment that exonerated the reaction--diffusion.}
\texttt{rd\_coral\_grow} rendered a transient grey band mid-growth that looked like a broken coral or a
shading bug. The decisive test was to strip it to a triple at \emph{identical} mechanics ---
\texttt{grow\_only}, \texttt{grow\_divide}, and \texttt{grow\_divide}$+$RD --- and compare
\emph{vertex trajectories}. The last two were byte-identical, so the reaction--diffusion is provably
pure colouring; \texttt{grow\_only} was smooth, so the roughness came from \emph{division under growth},
at scale. A per-frame trace nailed it: the cell count climbs in flat $+100$/frame steps (all cells
doubling together, a wave capped by \texttt{max\_div}) and the face-referenced roughness spikes to
$0.38$ (worst $1.7$) across the whole ramp. What the eye reported as ``grey'' was a
\emph{synchronisation}, not a mesh error.
\item \textbf{Stochastic cell cycle: the biological noise the model had omitted.} A deterministic
volume-doubling cycle with identical birth volumes \emph{must} synchronise under uniform growth; real
tissues don't, because cell-cycle times vary. Adding that variability --- a random initial phase at
seed (\texttt{vseed\_cv}) and a Gaussian per-division threshold (\texttt{cycle\_cv}) --- turns the
lockstep wave into a smooth ramp and drops the mid-ramp roughness $0.38\!\to\!0.004$, letting the coral
scale to $500\!\to\!1994$ cells cleanly. One parameter --- but only after the controlled experiment
localised the cause; guessing at shading or T1 tuning would have chased the symptom.
\item \textbf{Honest metrics matter (again).} The first roughness number was inflated by unused
reservoir vertices sitting at the origin. Restricting the metric to \emph{face-referenced} vertices ---
those that actually tile the shell --- gave the true $0.38\!\to\!0.004$ and matched the rendered strip.
Same lesson as the bow-tie faces below: measure the geometric thing you mean.
\end{itemize}

\paragraph{Stage 2 (2026-07-21) --- topology operators, and where the cheap division breaks.}
\begin{itemize}\itemsep2pt
\item \textbf{Revert-on-invalid is the right contract for topology.} Rather than reason about
every degenerate case up front, apply the flip/split, validate the touched faces, and roll back if
any is degenerate. T1 then \emph{never} tangles the mesh --- a refused flip is just a no-op. This
turned an error-prone operation into a safe one and is why the self-test is clean.
\item \textbf{Through-vertex division is a $1$--$2$ round operator, not a proliferation engine.}
Splitting a face by a septum between two of its \emph{existing} vertices adds no new vertices, so
each split roughly halves a cell using the same corners; after a couple of rounds cells become
triangles and further splits are refused. Worse, under \emph{uniform growth in a pinned patch} the
target area outruns the available space, every cell is driven below target, and the sliver
daughters collapse --- a run with $231$ forced divisions ended as a red tangle
($79$--$157$ invalid faces). The clean result above deliberately does \emph{one} desynchronised
round with no global growth. The proper fix (tyssue's own) is \emph{edge-midpoint} division:
bisect two opposite edges, insert new vertices (shared with the neighbours across those edges), and
run the septum between them --- well-shaped daughters, sustainable for many rounds. That needs a
vertex buffer with occupancy and neighbour-edge bookkeeping; it is the next increment.
\item \textbf{T1-mediated fluidisation (Stage~3) --- first not sound, then fixed by a \emph{bounded}
Euler step.} The naive attempt failed its own bar: gentle self-propulsion fired $0$ T1s, and the
drive needed to collapse edges made the floppy ($p_0\!=\!4.1$) sheet \emph{diverge} under a
fixed-step relaxation (positions overflow, $\sim\!85$ invalid faces) \emph{before} a T1 could
fluidise it --- the instability was in the \emph{integrator}, not the reconnection (near the
transition the shape energy is nearly flat, so an explicit step lets vertices cross). The sound
fix is a \emph{bounded overdamped Euler} step: cap every per-substep move at $15\%$ of the mean
edge length with \texttt{torch.clamp} --- so no vertex can cross a cell in one step, and the step
stays \emph{differentiable} (no \texttt{scipy} solver; the rollout must remain differentiable for
inverse modelling). With that, the sheet stays stable and T1s fire.
\item \textbf{A second, subtler bug the naked eye caught: bow-tie faces (a false ``valid'').} The
first ``sound'' dichotomy ($23\times$) was \emph{wrong}. The rendered strips showed black star-shaped
gaps where cells overlapped: the T1 flip was producing \emph{self-intersecting} (bow-tie) faces,
and my validity guard --- positive shoelace area, no repeated vertices --- \emph{passed} them,
because a bow-tie has positive shoelace area. So ``$0$ invalid faces'' was a false pass and the
big T1 counts were mostly bad flips. Two fixes: (i) a proper \emph{simple-polygon} test (no two
non-adjacent edges cross) plus a \emph{non-overlap} test on the four faces a flip touches, added to
the accept criterion; and (ii) the flip now \emph{searches} the insertion side (which side of its
anchor the gained vertex goes) $\times$ the mirror sign, and commits only a genuinely simple,
non-overlapping, edge-sharing result. With the honest check the corrected dichotomy is:
\begin{center}\small
\begin{tabular}{ccccc}
\toprule
$v_0$ & $l_{\text{th}}$ & solid T1 & fluid T1 & ratio \\
\midrule
0.06 & 0.25 & 96 & 1413 & $\mathbf{14.7\times}$ \\
\bottomrule
\end{tabular}
\end{center}
The fluid still performs $\sim\!15\times$ the T1s of the caged solid (fluid MSD $0.64$ vs solid
$0.10$ --- cage-breaking), now with a \emph{genuinely} valid, non-overlapping tiling at every frame
(the strips tile cleanly). The lesson is the one the whole document keeps re-learning: a numeric
invariant (area $>0$) is not the same as the geometric one you meant (a simple, embedded tiling) ---
and a rendered movie is a sharper validity check than a scalar. The corrected numbers are smaller
than the bogus ones, but the physics --- fluid flows, solid jams --- survives.
\end{itemize}

\paragraph{Stage 1 (2026-07-21) --- force balance is one line, and it just works.}
\begin{itemize}\itemsep2pt
\item The inner-relaxation-as-velocity trick (\S3) reaches force balance with no stability
tuning: $\Delta t$ is a free label, \texttt{eta}$=0.08$, $20$ inner steps/frame. Contrast the
SPV prototype, whose entire ``noisy pattern'' saga was a \emph{timestep} artefact. Moving the
DOF from cell centres to mesh vertices \emph{and} minimising to residual removed that whole
failure mode.
\item \textbf{The transition sits where the topology puts it.} With a fixed honeycomb (mostly
hexagons) the residual energy already collapses across $p_0\!\approx\!3.72$--$3.81$: the
single-cell floor for a regular hexagon is $p_0^\ast(\text{hex})=3.722$, so $\langle q\rangle$
cannot fall below $\approx\!3.72$ no matter how small $p_0$ is (the \texttt{solid} cases pin
there). The canonical $3.81$ is the \emph{disordered}-polygon value; sharpening the transition
to exactly $3.81$ needs T1 rearrangements to populate the pentagon/heptagon classes --- the next
operator. So the transition is not a fixed number but a readout of the polygon-class
distribution, which is a topology property.
\item \textbf{The engine calls \texttt{forward()} under \texttt{no\_grad}.} The autograd force
needs an explicit \texttt{with torch.enable\_grad()} (as the SPV op also did) --- otherwise
\texttt{torch.autograd.grad} raises ``does not require grad''. Cost one run to rediscover.
\item \textbf{Pin the border or the patch curls.} A free (unpinned) bounded sheet contracts
under perimeter tension and the boundary ring buckles inward, contaminating the bulk statistics.
Pinning vertices of degree $<3$ (those on the ragged Voronoi border) holds the frame; bulk stats
exclude any face touching a pinned vertex.
\end{itemize}

\paragraph{[next] morphogen-driven tubes and branching.}
The structural increments that were ``next'' in the earlier notes are now \emph{done}: the genuine
two-level hierarchy (a \texttt{cell} set carrying \texttt{chem}, fed by a \texttt{cell\_geometry}
Aggregate), edge-midpoint division with a decidable long-axis angle, and the lift to the closed 3D
monolayer (per-cell volume elasticity $+$ surface T1) all exist and are exercised above; the
reaction--diffusion runs on the cell set and rides division; and the division wave is tamed by a
stochastic cell cycle. What remains are the morphogenesis \emph{targets} from Turing\_vertex, and they
reuse every operator above unchanged --- only the coupling and the fluidity are new dials.
(i)~\emph{Morphogen-driven growth}: \texttt{morphogen\_growth\_3d} grows a cell's target volume where
the activator $a$ is high, so a localised Turing patch on a \emph{fluid} ($p_0\!>\!3.81$) shell, where
T1 lets the tissue flow and its cells divide (handing the activator to their daughters), should
\emph{bud and extend into a tube} rather than jam into a spike (Turing\_vertex Fig.~5). (ii)~Lowering
the reaction--diffusion feed drives that patch to \emph{branch} (Fig.~6). The exploratory \texttt{tube}
runner --- the same schedule with the morphogen$\to$growth coupling switched on and $p_0$ raised into
the fluid regime --- is the current front; stabilising the bud against the fluid shell's tendency to
balloon is the open problem.

\end{document}
